\section{Conclusion}

We have developed an internal language for representable
semilattice-enriched multicategories.  Its distributive rules absorb
multilinearity into the syntax, at the cost of non-unique derivations.  We
identify derivations by permutation equivalence and define substitution on
them; its result is derivation-independent up to the semilattice equations and
satisfies the multicategory laws.  The quotient calculus presents the free
representable $\catname{SLat}$-multicategory on a multigraph.

In the cartesian instance, $\join$ collects the representatives belonging to
an e-class, and the equations of an e-graph presentation become equations or
rewrite rules in the calculus.  Equality saturation can therefore be viewed
as term rewriting modulo the structural equations of the type theory.

The definitional extension \textbf{let} records acyclic sharing without
changing the category.  The recursive \textbf{letrec} captures cyclic
e-graphs; its recursive joins remain protected because trace need not preserve
joins.  This calculus presents the free traced cartesian representable
$\catname{SLat}$-multicategory, equivalently its traced cartesian monoidal
form.

This gives a syntactic counterpart to the string-diagram semantics of Tiurin
et al.~\cite{tiurin2025equivalencehypergraphsdporewriting}, and a foundation
for further type-theoretic accounts of e-graphs and equality saturation.

\paragraph{Future work.}
Finally, we speculate that this work may lay some of the mathematical groundwork 
for programming languages with e-graphs built in as a primitive concept.
Just as the simply typed lambda calculus underlies the design of functional 
programming languages, a type theory of the kind developed here could inform the 
design of a language where equality saturation is not an external optimization 
pass but an intrinsic part of the language's evaluation semantics---one where 
terms natively denote equivalence classes of programs rather than individual 
programs.
Adding functions to our type theory is routine as they only require a corresponding quotioenting for the universal property of a closed multicategory.
We leave the development of such a language as an ambitious direction for future 
work.
